\documentclass[11pt]{article}

\usepackage[a4paper,margin=2.2cm]{geometry}
\usepackage{amsmath}
\usepackage{booktabs}
\usepackage{graphicx}
\usepackage[hidelinks]{hyperref}
\usepackage{seqsplit}
\usepackage{url}

\hypersetup{
  pdftitle={Supplementary Material: Confidence-Shielded Evaluator-Conditioned Fusion for Risk-Controlled Infrared-Visible Object Detection},
  pdfauthor={Jianyong Zhao, Zigao Wu, Mengfei Ma},
  pdfsubject={Supplementary material for a research paper submitted to Infrared Physics \& Technology}
}

\newcommand{\code}[1]{\texttt{#1}}

\title{Supplementary Material for\\Confidence-Shielded
Evaluator-Conditioned Fusion for Risk-Controlled Infrared--Visible Object
Detection}
\author{Jianyong Zhao, Zigao Wu, Mengfei Ma}
\date{}

\begin{document}
\maketitle
\sloppy

\section{Data provenance and split control}

LLVIP uses a fixed 1,350/150/300 train/validation/test partition. The manifest
contains the complete stem lists and SHA-256 digests. All three pairwise stem
intersections are empty. The 300 evaluated images are the official test images;
they are not used for checkpoint selection.

\begin{table}[h]
\centering
\small
\caption{Partition sizes and annotation counts.}
\begin{tabular}{@{}lrrrr@{}}
\toprule
Dataset & Training & Validation & Test & Test groups \\
\midrule
LLVIP & 1,350 & 150 & 300 & 7 scenes \\
M3FD & 1,600 & 400 & 200 & 2 frame groups \\
DroneVehicle & public train & public validation & 300 & 3 frame groups \\
\bottomrule
\end{tabular}
\end{table}

Using the same training-run-cluster and group bootstrap for the primary
CS-ECF-versus-visible contrast gives a dataset-balanced AP50 gain of 0.0133 and
a 95\% interval of \([0.0026,0.0261]\). This clustered interval is the primary
uncertainty statement; treating the two DroneVehicle checkpoints as independent
would understate their shared training-run dependence.

\begin{table}[h]
\centering
\small
\caption{LLVIP annotation counts in the fixed split and pairwise overlap.}
\begin{tabular}{@{}lrl@{}}
\toprule
Split & Pedestrian boxes & Pairwise stem overlap \\
\midrule
Training & 4,479 & train--validation: 0 \\
Validation & 419 & validation--test: 0 \\
Test & 744 & train--test: 0 \\
\bottomrule
\end{tabular}
\end{table}

M3FD frames are partitioned into consecutive 100-frame groups. The initial
grouping seed is 2028 and the extension seed is 2029; the two held-out test
groups do not enter training or validation. DroneVehicle uses three consecutive
100-frame groups selected from the official test archive with seed 2037. The
source archive has 4,430,343,899 bytes and SHA-256
\texttt{\seqsplit{94BF47E493E7A57FD5D900439F614F339E7B6C7181E4D781BCFEBAEB963FFD8E}}.

The YOLO11x-OBB best and last checkpoints share one 150-epoch training run
(batch size 32, 640-pixel input, seed 0). They are treated as one training-run
cluster in every aggregate bootstrap.

\section{Detection evaluation}

All predictions are generated at a confidence floor of 0.001 and a 640-pixel
inference size. YOLO predictions use NMS at IoU 0.7. For each image, class and
IoU threshold from 0.50 to 0.95 in increments of 0.05, predictions are assigned
to ground truth by one-to-one maximum-cardinality matching over all valid
prediction--annotation pairs. The resulting true positives and confidence
scores are pooled across the evaluated images, and AP is computed on the
101-point interpolated precision--recall curve with Ultralytics 8.4.189.
AP50 averages class AP at IoU 0.50; AP50--95 additionally averages the ten IoU
thresholds.

\section{Complete CS-ECF and Always CDF results}

The tables below reconstruct the frozen policy from released per-image
detections. ``Selected'' gives the number of images assigned to visible, CDF,
PIAFusion and SPGFusion. The group intervals are percentile bootstraps over
scenes or consecutive frame groups.

\begin{table}[ht]
\centering
\small
\caption{Per-checkpoint paired comparison on LLVIP. Delta is CS-ECF minus
Always CDF; the last column is Visible/CDF/PIAFusion/SPGFusion selections.}
\label{tab:llvip-cdf}
\begin{tabular}{@{}lrrrrr@{}}
\toprule
Evaluator & CDF & CS-ECF & Delta & 95\% group CI & V/C/P/S \\
\midrule
YOLO11n & .8920 & .8845 & \(-.0076\) & \([-.024,+.012]\) & 276/1/9/14 \\
Agnostic & .8629 & .8611 & \(-.0018\) & \([-.022,+.011]\) & 286/2/6/6 \\
YOLO26n s0 & .8200 & .8088 & \(-.0112\) & \([-.034,+.009]\) & 288/1/3/8 \\
RT-DETR-l & .8100 & .8136 & \(+.0036\) & \([-.015,+.018]\) & 293/0/0/7 \\
YOLO26n s1 & .8048 & .7911 & \(-.0138\) & \([-.026,+.002]\) & 265/1/4/30 \\
YOLO26n s2 & .8183 & .8087 & \(-.0096\) & \([-.023,+.012]\) & 282/0/0/18 \\
YOLOv8n s0 & .7847 & .7761 & \(-.0086\) & \([-.025,+.011]\) & 282/1/0/17 \\
YOLOv8n s1 & .8087 & .8119 & \(+.0032\) & \([-.015,+.024]\) & 290/0/0/10 \\
\bottomrule
\end{tabular}
\end{table}

\begin{table}[ht]
\centering
\small
\caption{Per-checkpoint paired comparison on M3FD and DroneVehicle.}
\label{tab:external-cdf}
\begin{tabular}{@{}llrrrr@{}}
\toprule
Dataset & Evaluator & CDF & CS-ECF & Delta & 95\% group CI \\
\midrule
M3FD & YOLO26n & .3705 & .3757 & \(+.0053\) & \([-.0002,+.0114]\) \\
M3FD & YOLOv8n & .3731 & .3972 & \(+.0241\) & \([+.0087,+.0321]\) \\
DroneVehicle & best & .6624 & .6814 & \(+.0190\) & \([-.0013,+.0301]\) \\
DroneVehicle & last & .6524 & .6767 & \(+.0243\) & \([-.0067,+.0464]\) \\
\bottomrule
\end{tabular}
\end{table}

\begin{table}[h]
\centering
\small
\caption{Paired CS-ECF minus Always CDF summary.}
\begin{tabular}{@{}lrc@{}}
\toprule
Scope & Difference & 95\% interval \\
\midrule
LLVIP & \(-0.0057\) & \([-0.0189,+0.0156]\) \\
M3FD & \(+0.0147\) & \([+0.0100,+0.0159]\) \\
DroneVehicle & \(+0.0217\) & \([-0.0007,+0.0738]\) \\
Dataset-balanced & \(+0.0102\) & \([-0.00057,+0.02980]\) \\
\bottomrule
\end{tabular}
\end{table}

The combined interval resamples scene or frame groups and training-run clusters
within each dataset, then assigns equal weight to the three datasets. There are
eight run clusters on LLVIP, two on M3FD and one on DroneVehicle. The
DroneVehicle component therefore quantifies test-group uncertainty but cannot
estimate variability between independent DroneVehicle training runs. The lower
endpoint of the combined interval is below zero, so the direct CDF comparison
does not establish a separated average improvement.

\section{External evaluator stress tests}

Two COCO-pretrained detectors, YOLOv9t and YOLOv10s, were used without LLVIP
fine-tuning. They were not used in policy construction, threshold selection, or
the primary evaluation. Person-class detections were retained before AP and
selection scoring. The 1.5 margin and top-five score were unchanged.

A third stress test used an LLVIP-fine-tuned YOLOv9t. The checkpoint was
initialized from COCO YOLOv9t and trained on the 1,350 visible training images
for a pre-specified 40 epochs. Training used AdamW, initial learning rate
0.002, three-epoch warm-up, cosine decay, weight decay \(5\times10^{-4}\), batch
size 16, 640-pixel inputs, random horizontal flipping, and
brightness/contrast jitter. A reproducible PIL-based loader calling the
Ultralytics detection loss was used. The final checkpoint after epoch 40 was
frozen before the LLVIP test images were scored; no epoch was selected using
test results. Its SHA-256 is
\texttt{\seqsplit{AE0A44359C489DEEAE3F7443173588B9F68E264B6BE1EA455E961EB2DED0797D}}.
The checkpoint reached mAP50 0.8456 on the 150-image validation set and 0.6666
on the 300-image test visible input.

\begin{table}[ht]
\centering
\small
\caption{External evaluator stress tests on the same 300 LLVIP test pairs.}
\begin{tabular}{@{}lrrrr@{}}
\toprule
Model & Visible & Always CDF & No shield & CS-ECF \\
\midrule
YOLOv9t, zero-shot & .7698 & .7888 & .8216 & .7948 \\
YOLOv10s, zero-shot & .8344 & .8617 & .8699 & .8557 \\
YOLOv9t, LLVIP-trained & .6666 & .6598 & .6405 & .6558 \\
\bottomrule
\end{tabular}
\end{table}

\begin{table}[ht]
\centering
\scriptsize
\setlength{\tabcolsep}{3pt}
\caption{Paired confidence intervals and scene-tail behavior for the external
evaluators. The final column reports the worst scene change for CS-ECF /
no-shield.}
\begin{tabular}{@{}lccc@{}}
\toprule
Model & CS-ECF \(-\) visible & CS-ECF \(-\) CDF & Worst (CS/NS) \\
\midrule
YOLOv9t, zero-shot & \(+.0250~[+.0053,+.0523]\) &
  \(+.0060~[-.0277,+.0374]\) & \(-.0167 / -.0011\) \\
YOLOv10s, zero-shot & \(+.0213~[+.0041,+.0439]\) &
  \(-.0059~[-.0266,+.0180]\) & \(-.0038 / -.0337\) \\
YOLOv9t, LLVIP-trained & \(-.0108~[-.0342,+.0050]\) &
  \(-.0040~[-.0206,+.0108]\) & \(-.0724 / -.1208\) \\
\bottomrule
\end{tabular}
\end{table}

The shield selected visible on 233/300 zero-shot YOLOv9t, 245/300 YOLOv10s,
and 230/300 LLVIP-trained YOLOv9t images. The two zero-shot evaluators support
a positive difference from visible but not superiority over Always CDF. The
fine-tuned YOLOv9t is a negative portability case: CS-ECF was below both
visible and Always CDF, with scene-group intervals crossing zero. The shield
improved the unshielded selector by 0.0153 AP50 (group interval
\([+0.0072,+0.0267]\)) and reduced the worst scene loss from 0.1208 to 0.0724,
but did not remove the loss. Tail behavior therefore remains evaluator
dependent rather than universally controlled.

\section{Validation choice and exploratory alternatives}

The policy was fixed before the confirmatory comparisons. On validation, the
four-candidate top-five selector with a 1.5 visible margin produced gains of
0.001850 and \(-0.000070\) for the two LLVIP evaluators and 0.001518 and
0.000083 for the two M3FD evaluators. Its two-dataset balanced validation
average was 0.000845. The candidate set was retained to preserve the visible,
region-selected, learned, and published fusion families; the margin was fixed
at 1.5.

Two alternative policies were evaluated after the frozen comparison. They are
exploratory and are not used to support the main claim. A top-10 selector over
visible/CDF/PIAFusion with margin 1.3 gave a 0.0039 balanced gain but a
hierarchical interval of \([-0.0029,+0.0127]\). A mean-confidence selector over
visible/SPGFusion with margin 1.5 gave a \(-0.0034\) balanced gain and
\([-0.0085,+0.0013]\). Both alternatives fail to establish a stable advantage
across the 10 LLVIP/M3FD checkpoints.

\section{Thermal controls}

\begin{table}[h]
\centering
\small
\caption{Spearman association between switching and thermal or visible
features. These secondary association tests are descriptive and are not
multiplicity-adjusted.}
\begin{tabular}{@{}lrrr@{}}
\toprule
Dataset & Thermal top-hat \(\rho\) (\(p\)) & Visible luma \(\rho\) (\(p\)) &
CDF-mask \(\rho\) (\(p\)) \\
\midrule
LLVIP & \(+0.151\) (0.009) & \(-0.104\) (0.071) & \(-0.029\) (0.614) \\
M3FD & \(+0.320\) (\(<10^{-5}\)) & \(-0.066\) (0.356) & \(-0.332\) (\(<10^{-5}\)) \\
DroneVehicle & \(-0.403\) (\(<10^{-5}\)) & \(-0.441\) (\(<10^{-5}\)) & \(+0.064\) (0.269) \\
\bottomrule
\end{tabular}
\end{table}

\noindent In the fixed-input CDF ablation on DroneVehicle, the thermal mask
changed AP50 from 0.6620 to 0.6624 for the best checkpoint and from 0.6512 to
0.6524 for the last checkpoint, relative to global detail injection.

The association reverses between datasets. Thermal contrast is therefore not
treated as a sufficient proxy for fusion utility, and no universal
thermal-switching rule is claimed.

\section{Reproduction}

The public repository contains the selector, candidate-generation and detector
scoring code, per-image score arrays, split manifests, source data for every
figure and the scripts that reproduce the intervals. The direct baseline
comparison is generated by:

\begin{verbatim}
python scripts/analyze_csecf_vs_cdf.py
\end{verbatim}

This command validates each checkpoint AP before writing
\path{results/selective_fusion/csecf_vs_cdf_paired.json}. Split and candidate
manifests retain their seeds and selected groups. The independent evaluator
checkpoints were trained or loaded by
\path{scripts/train_independent_manual.py}, scored with
\path{scripts/score_clean_evaluators.py}, and summarized with
\path{scripts/analyze_independent_evaluators.py}. The LLVIP-trained YOLOv9t
result is stored in
\path{results/independent_evaluators/domain_yolov9t_v3/}.

\end{document}
